\section{Conclusions}
%
We proposed a fairness measure that accomplishes two important desiderata.
First, it remedies the main conceptual shortcomings of demographic parity as a
fairness notion. Second, it is fully aligned with the central goal of supervised
machine learning, that is, to build higher accuracy classifiers.
In light of our results, we draw several conclusions aimed to help interpret and
apply our framework effectively.

\paragraph{Choose reliable target variables.}

Our notion requires access to observed outcomes such as default rates in the
loan setting. This is precisely the same requirement that supervised learning
generally has. The broad success of supervised learning demonstrates that this
requirement is met in many important applications. That said, having access
to reliable ``labeled data'' is not always possible. Moreover, the measurement
of the target variable might in itself be unreliable or biased. Domain-specific
scrutiny is required in defining and collecting a reliable target variable.

\paragraph{Measuring unfairness, rather than proving fairness.}

Due to the limitations we described, satisfying our notion (or any other
oblivious measure) should not be considered a conclusive \emph{proof of
fairness}. Similarly, violations of our condition are not meant to be a proof of
unfairness. Rather we envision our framework as providing a reasonable way of
discovering and measuring potential concerns that require further scrutiny. We
believe that resolving fairness concerns is ultimately impossible without
substantial domain-specific investigation. This realization echoes earlier
findings in ``Fairness through Awareness''~\cite{DworkHPRZ12} describing the
task-specific nature of fairness.

\paragraph{Incentives.}

Requiring equalized odds creates an incentive structure for the entity
building the predictor that aligns well with achieving fairness.
Achieving better prediction with equalized odds requires collecting
features that more directly capture the target $Y$, unrelated to its
correlation with the protected attribute. Deriving an equalized
odds predictor from a score involves considering the pointwise minimum
ROC curve among different protected groups, encouraging constructing
of predictors that are accurate in all groups, e.g.,~by collecting data
appropriately or basing prediction on features predictive in all
groups.

\paragraph{When to use our post-processing step.}
An important feature of our notion is that it can be achieved via a simple and
efficient post-processing step. In fact, this step requires only aggregate
information about the data and therefore could even be carried out in a
privacy-preserving manner (formally, via Differential Privacy). In contrast,
many other approaches require changing a usually complex machine learning
training pipeline, or require access to raw data. Despite its simplicity, our
post-processing step exhibits a strong optimality principle. If the underlying
score was close to optimal, then the derived predictor will be close to optimal
among all predictors satisfying our definition. However, this does not mean that
the predictor is necessarily good in an absolute sense. It also does not mean
that the loss compared to the original predictor is always small.  An
alternative to using our post-processing step is always to invest in better
features and more data. Only when this is no longer an option, should our
post-processing step be applied.

\paragraph{Predictive affirmative action.}
In some situations, including Scenario II in
Section~\ref{sec:oblivious}, the equalized odds predictor can be thought of as
introducing some sort of affirmative action: the optimally predictive
score~$R^*$ is shifted based on~$A$. This shift compensates for the fact that,
due to uncertainty, the score is in a sense more biased then the target label
(roughly, $R^*$ is more correlated with~$A$ then~$Y$ is correlated with~$A$).
Informally speaking, our approach transfers the \emph{burden of uncertainty}
from the protected class to the decision maker. We believe this is a reasonable
proposal, since it incentivizes the decision maker to invest additional
resources toward building a better model.
